% C.3be. Complete draft; independent full proof and source audits requested.
\paragraph{C.3be. Rational top Chow groups over
\(\overline{\mathbb F}_p\) in every local dimension (Proved).}
Let \(k=\overline{\mathbb F}_p\), and let \(R\) be the strict
henselization at a closed \(k\)-point of a normal local ring
essentially of finite type over \(k\). Suppose that
\(N=\dim R\geq2\), and put
\[
 S=\operatorname{Spec}R,\qquad
 X=S\setminus\{\mathfrak m_R\}.
\]
Then
\begin{equation}\label{eq:C3be-top-vanishing}
       \mathrm{CH}_1(S)_{\mathbb Q}=0,
       \qquad \mathrm{CH}^{N-1}(X)_{\mathbb Q}=0,
       \qquad B^{N-1}(X)_{\mathbb Q}=0.
\end{equation}
Here the dimension-indexed Chow groups use the dimension function
\(\delta(\mathfrak p)=\dim(R/\mathfrak p)\) on \(S\), and its
restriction on \(X\). In particular
\(\mathrm{CH}_1(X)=\mathrm{CH}^{N-1}(X)\); this is not an
identification obtained by subtracting one from the Krull dimension
of the punctured spectrum. No regularity assumption on \(X\), and
no resolution theorem in dimension greater than two, is required.
All positive characteristics, including \(2\) and \(3\), are included.

\begin{proof}
\emph{1. Normalized surface quotients remain algebraizable.}
We first prove the finite-model assertion needed to apply C.3bd.
Let \(\mathfrak q\subset R\) be a prime with
\(\dim(R/\mathfrak q)=2\), and set \(B=R/\mathfrak q\).
Write \(R=A^{\mathrm{sh}}\) for the given local model. Since
\(k\) is algebraically closed, strict henselization and
henselization agree here. The filtered pointed-\'etale construction
\citep[Lemmas 10.155.1--2]{frC3beStacksHenselConstruction}
allows us to choose an essentially \(\acute{e}\)tale local
\(A\)-algebra \(A'\), with residue field \(k\) and
\((A')^{\mathrm{sh}}=R\), containing a finite set of generators
of \(\mathfrak q\). Let \(I\subset A'\) be the ideal generated
by these elements. Thus \(IR=\mathfrak q\).

The map \(A'\to R\) is faithfully flat. Hence
\[
                 I=IR\cap A'=\mathfrak q\cap A'
\]
is prime. Put \(B_0=A'/I\). Henselization commutes with
quotients, so
\begin{equation}\label{eq:C3be-surface-model}
                         B=B_0^{\mathrm{sh}}
\end{equation}
by \citet[Lemma 10.156.4, Tag 05WS]
{frC3beStacksQuotientHenselization}.
The dimension is preserved by this local ind-\'etale map, so
\(B_0\) is a two-dimensional local domain essentially of finite
type over \(k\).

Let \(\overline B_0\) and \(\overline B\) denote the
normalizations. Excellence makes both normalizations finite.
By \citet[Lemma 37.19.4, Tag 0CBM]
{frC3beStacksNormalizationHenselization},
\[
                 \overline B
                    =\overline B_0\otimes_{B_0}B.
\]
Since \(B\) is a henselian domain, its integral normalization
\(\overline B\) is a henselian local domain
\citep[Lemma 15.106.23, Tag 092X]
{frC3beStacksIntegralHenselian}.
It has residue field \(k\): the residue extension is finite,
and \(k\) is algebraically closed.
Choose the point \(v\) of \(\operatorname{Spec}\overline B_0\)
under this unique closed point. The finite local-factor
description of strict henselization gives
\begin{equation}\label{eq:C3be-normalized-model}
              \overline B
                 \simeq(\overline B_{0,v})^{\mathrm{sh}};
\end{equation}
we use \citet[Lemma 10.156.3, Tag 05WR]
{frC3beStacksFiniteHenselization}.
Thus \(\overline B\) is precisely a normal algebraizable
strict-local surface of the type covered by C.3bd. This argument
does not merely assert that a finite normalization is henselian;
it supplies the required normal local model of finite type.

\emph{2. The rational divisor group of each normalized surface vanishes.}
The puncture of the normal surface
\(\operatorname{Spec}\overline B\) is regular.
C.3bd in dimension two therefore gives
\[
 \mathrm{CH}^{1}\!left(
    \operatorname{Spec}\overline B
       \setminus\{\mathfrak m_{\overline B}\}\right)_{\mathbb Q}=0.
\]
The omitted closed point has \(\delta\)-dimension zero.
Chow localization
\citep[Lemma 42.19.3, Tag 02RX]{frC3beStacksChowLocalization}
therefore identifies this group with
\(\mathrm{CH}_1(\operatorname{Spec}\overline B)_{\mathbb Q}\).
In particular every prime divisor on this normalized local
surface is torsion in its Chow group. We claim no uniform
annihilating integer for all such divisors.

\emph{3. Every curve is killed on a surface containing it.}
Let \(\mathfrak p\subset R\) satisfy
\(\dim(R/\mathfrak p)=1\). Since \(R\) is a catenary local
domain of dimension \(N\), a saturated chain through
\(\mathfrak p\) contains a prime
\[
                   \mathfrak q\subset\mathfrak p,
                    \qquad\dim(R/\mathfrak q)=2.
\]
For \(N=2\) we take \(\mathfrak q=(0)\).
Equivalently, catenarity gives
\(\operatorname{ht}(\mathfrak p)=N-1\), and we take the
preceding prime in a saturated chain. It has height \(N-2\).

Use the normalization \(\overline B\) of
\(B=R/\mathfrak q\) from Step~1. By lying over there is a
prime \(\mathfrak P\subset\overline B\) above
\(\mathfrak p/\mathfrak q\). It is neither zero nor maximal,
so it has height one in the normal local surface
\(\operatorname{Spec}\overline B\). The finite morphism
\[
       g:\operatorname{Spec}\overline B
                 \longrightarrow\operatorname{Spec}R
\]
is the normalization followed by the closed immersion of the
surface quotient. Proper pushforward gives
\begin{equation}\label{eq:C3be-curve-push}
       g_*[V(\mathfrak P)]
           =[\kappa(\mathfrak P):\kappa(\mathfrak p)]
                                      [V(\mathfrak p)].
\end{equation}
The coefficient is a positive finite integer. Proper pushforward
preserves rational equivalence over the universally catenary
base \((S,\delta)\)
\citep[Lemma 42.20.3, Tag 02S2]{frC3beStacksProperPush}.
The left side of \eqref{eq:C3be-curve-push} is zero with rational
coefficients by Step~2. Consequently
\([V(\mathfrak p)]=0\) in \(\mathrm{CH}_1(S)_{\mathbb Q}\).
These classes generate the group, proving its vanishing.

Finally the closed point of \(S\) carries no
\(\delta\)-dimension-one cycles. The same Chow localization
sequence gives
\[
           \mathrm{CH}_1(S)\xrightarrow{\sim}
               \mathrm{CH}_1(X)=\mathrm{CH}^{N-1}(X).
\]
The final ordinary quotient vanishes as well. The proof has used
only normal surface resolutions through C.3bd; the ambient
dimension \(N\) enters solely in the catenary choice of a
surface through the given curve.
\end{proof}
